% !TeX program = pdflatex
\documentclass[9pt,a4paper]{extarticle}
\usepackage[spanish,es-nodecimaldot,es-noshorthands]{babel}
\usepackage[T1]{fontenc}
\usepackage[utf8]{inputenc}
\usepackage{lmodern}
\renewcommand{\familydefault}{\sfdefault}
\usepackage[a4paper,left=0.92cm,right=0.92cm,top=0.7cm,bottom=0.76cm,headheight=9pt]{geometry}
\usepackage{graphicx}
\usepackage{xcolor}
\usepackage{tikz}
\usetikzlibrary{positioning,arrows.meta}
\usepackage{tcolorbox}
\tcbuselibrary{skins}
\usepackage{booktabs}
\usepackage{tabularx}
\usepackage{array}
\usepackage{enumitem}
\usepackage{fancyhdr}
\usepackage{hyperref}
\usepackage{url}
\usepackage{microtype}

\definecolor{cchenpurple}{HTML}{483888}
\definecolor{cchenlime}{HTML}{DDDC40}
\definecolor{textgray}{HTML}{333333}
\definecolor{mutedgray}{HTML}{6F6F6F}
\definecolor{lightgray}{HTML}{F3F3F5}
\definecolor{linegray}{HTML}{D9D9DF}
\definecolor{okgreen}{HTML}{2E7D32}
\definecolor{warnamber}{HTML}{F9A825}
\definecolor{badred}{HTML}{C62828}

\hypersetup{colorlinks=true, linkcolor=cchenpurple, urlcolor=cchenpurple,
  pdftitle={Ficha de fuente de información - PatentsView / USPTO},
  pdfauthor={CCHEN 360}}
\urlstyle{same}
\pagestyle{fancy}
\fancyhf{}
\fancyfoot[L]{\scriptsize\color{mutedgray}Observatorio CCHEN 360 -- ficha de fuente de información}
\fancyfoot[R]{\scriptsize\color{mutedgray}\thepage}
\renewcommand{\headrulewidth}{0pt}
\renewcommand{\footrulewidth}{0.25pt}
\setlength{\parindent}{0pt}
\setlength{\parskip}{1.2pt}
\hyphenpenalty=10000
\exhyphenpenalty=10000
\setlist[itemize]{leftmargin=*,topsep=1pt,itemsep=0.2pt,parsep=0pt}
\renewcommand{\arraystretch}{1.04}
\newcolumntype{Y}{>{\raggedright\arraybackslash}X}
\newcommand{\sectiontitle}[1]{\vspace{3pt}{\normalsize\bfseries\color{cchenpurple}#1}\par\vspace{1pt}}
\newtcolorbox{keybox}{enhanced,colback=cchenpurple!5,colframe=cchenpurple,boxrule=0.5pt,arc=2pt,left=4pt,right=4pt,top=3pt,bottom=3pt}
\newtcolorbox{metricbox}[2]{enhanced,colback=#1!7,colframe=#1,boxrule=0.8pt,arc=2pt,left=5pt,right=5pt,top=4pt,bottom=4pt}
\newtcolorbox{calloutbox}[1]{enhanced,colback=lightgray,colframe=linegray,boxrule=0.4pt,arc=2pt,left=4pt,right=4pt,top=3pt,bottom=3pt,title={#1},fonttitle=\bfseries\footnotesize\color{cchenpurple},attach title to upper={\par\vspace{0.5pt}}}
\newcommand{\dotpattern}{\begin{tikzpicture}[remember picture,overlay]\foreach \x in {0.2,0.55,...,11.2} {\foreach \y in {-0.1,0.22,0.54,0.86} {\fill[cchenpurple!55] (\x,\y) circle (0.55pt); \fill[cchenlime!85] ({\x+0.16},{\y+0.13}) circle (0.48pt);}}\end{tikzpicture}}

\begin{document}
\thispagestyle{fancy}

\begin{minipage}[t]{0.34\textwidth}
  \includegraphics[height=0.9cm]{../../../licitacion/assets/ministerio\_energia\_color.png}
\end{minipage}
\begin{minipage}[t]{0.36\textwidth}
  \vspace{0.08cm}\dotpattern
\end{minipage}
\begin{minipage}[t]{0.28\textwidth}
  \raggedleft\includegraphics[height=0.9cm]{../../../licitacion/assets/cchen360\_logo\_color.png}
\end{minipage}

\vspace{0.11cm}
{\color{cchenpurple}\rule{\textwidth}{1.5pt}}
\vspace{0.05cm}

\begin{center}
{\fontsize{16}{17.5}\selectfont\bfseries\color{cchenpurple}PatentsView / USPTO}\\[-1pt]
{\fontsize{9.1}{10.4}\selectfont\color{textgray}Ficha de fuente de información para datos relacionados con CCHEN}\\[-1pt]
{\fontsize{7.8}{9}\selectfont\color{mutedgray}Código: \texttt{patentsview\_uspto} \quad | \quad Categoría: Patentes}
\end{center}
\vspace{0.01cm}

\begin{keybox}
\textbf{\color{cchenpurple}Mensaje operativo.} Esta fuente se documenta solo para registros relacionados con CCHEN. No se descarga ni se usa el universo completo: el filtro opera por DOI, autor, afiliación, ROR/ORCID, alias institucional, resultados conocidos o semillas temáticas justificadas.
\end{keybox}

\vspace{0.04cm}
\begin{minipage}{0.24\textwidth}
\begin{metricbox}{okgreen}{Datos obtenidos}
{\large\bfseries Datos obtenidos}\par
{\Large\bfseries 0}\par
{\normalsize conteo registrado}
\end{metricbox}
\end{minipage}\hfill
\begin{minipage}{0.24\textwidth}
\begin{metricbox}{cchenpurple}{Frecuencia}
{\large\bfseries Frecuencia}\par
{\large\bfseries semestral}\par
{\normalsize refresco sugerido}
\end{metricbox}
\end{minipage}\hfill
\begin{minipage}{0.24\textwidth}
\begin{metricbox}{warnamber}{Tipo}
{\large\bfseries Tipo}\par
{\normalsize\bfseries Bloqueada por token}\par
{\normalsize implementación}
\end{metricbox}
\end{minipage}\hfill
\begin{minipage}{0.24\textwidth}
\begin{metricbox}{badred}{Decisión}
{\large\bfseries Decisión}\par
{\normalsize\bfseries Bloqueada por credencial}\par
{\normalsize criterio operativo}
\end{metricbox}
\end{minipage}

\vspace{0.04cm}
\begin{minipage}[t]{0.48\textwidth}\raggedright
\sectiontitle{1. Descripción de la fuente}
Interfaz de consulta de patentes USPTO. Permite buscar propiedad industrial, inventores y solicitantes cuando existe credencial de acceso.

\sectiontitle{2. Qué datos se descargaron}
\begin{calloutbox}{Tipología de datos}
Patentes USPTO, solicitantes e inventores
\end{calloutbox}
\textbf{Datos descargados para CCHEN.} Salida prevista para patentes USPTO CCHEN; actualmente sin registros por falta de credencial.

\sectiontitle{3. Cómo se filtra}
Aliases de solicitante/inventor CCHEN; requiere credencial PatentsView.

\sectiontitle{4. Dónde quedó guardado}
\begin{itemize}
\item \path{Data/Patents/cchen_patents_uspto.csv}
\end{itemize}

\sectiontitle{5. Estadística de descarga}
Registros obtenidos/evidencia: 0. Archivos locales generados o registrados: 1. Calidad registrada: 0.0. Última actualización: 2026-05-19T10:54:05. Estado de corrida: failed.

\end{minipage}\hfill
\begin{minipage}[t]{0.48\textwidth}\raggedright
\sectiontitle{6. Para qué sirve en el observatorio}
\begin{itemize}
  \item \textbf{Uso principal:} Serviria para propiedad industrial internacional cuando exista credencial; por ahora queda como brecha documentada.
  \item \textbf{Potencial:} Buscar patentes USPTO asociadas a aliases CCHEN; actualmente depende de credencial.
\end{itemize}

\sectiontitle{7. Estado operativo}
\begin{calloutbox}{Estado}
Registrada, pero bloqueada hasta configurar credencial/credencial.
\end{calloutbox}
Estado catalogo: implementada en registro operativo. Ultima corrida: failed; ultima actualizacion: sin fecha.

\sectiontitle{8. Debilidades y riesgos}
Registrada, pero la corrida queda bloqueada hasta configurar credencial PatentsView.

\sectiontitle{9. Decisión}
Mantener registrada como brecha operativa. Activar cuando exista credencial gratuita y separar claramente de INAPI local.

\end{minipage}

\vspace{0.06cm}
\begin{tcolorbox}[enhanced,colback=white,colframe=cchenpurple!55,boxrule=0.45pt,arc=2pt,left=6pt,right=6pt,top=5pt,bottom=5pt]
{\bfseries\color{cchenpurple}Método replicable:} tomar la fuente, filtrar solo datos relacionados con CCHEN, guardar el archivo local, revisar calidad y usar la evidencia para indicadores del observatorio.
\end{tcolorbox}

\vspace{0.02cm}
{\scriptsize\color{mutedgray}Sitio: https://patentsview.org; https://uspto.gov \quad | \quad Interfaz de datos: https://patentsview.org/apis/purpose; https://developer.uspto.gov/api-catalog}

\end{document}
